\documentclass[12pt,a4paper]{article}
\usepackage{../../../myconfig}

\begin{document}

% --- KHUNG TIÊU ĐỀ ĐẦU TRANG ---
\titlebox{Chủ đề: Oxyz}{Phương trình đường thẳng}{Oxyz: Phương trình đường thẳng}

% =========================================================================
% PHẦN 1: NHẬN BIẾT - VTCP & ĐIỂM THUỘC ĐƯỜNG THẲNG (CÂU 1 - CÂU 10)
% =========================================================================

\questionLinesManual{0}{1}{%
Trong không gian $Oxyz$, cho đường thẳng $d\colon \dfrac{x-2}{3}=\dfrac{y+5}{4}=\dfrac{z-2}{-1}$. Vectơ nào dưới đây là một vectơ chỉ phương của $d$?
}{%
\begin{tasks}(4)
    \task $\vec{u}_1=(3;4;-1)$
    \task $\vec{u}_2=(2;-5;2)$
    \task $\vec{u}_3=(2;5;-2)$
    \task $\vec{u}_4=(3;4;1)$
\end{tasks}
}

\questionLinesManual{0}{2}{%
Trong không gian $Oxyz$, cho đường thẳng $d\colon \begin{cases} x=2-t \\ y=1+2t \\ z=3+t \end{cases}$. Vectơ nào dưới đây là một vectơ chỉ phương của $d$?
}{%
\begin{tasks}(4)
    \task $\vec{u}_1=(-1;2;3)$
    \task $\vec{u}_2=(2;1;3)$
    \task $\vec{u}_3=(-1;2;1)$
    \task $\vec{u}_4=(2;1;1)$
\end{tasks}
}

\questionLinesManual{0}{3}{%
Trong không gian $Oxyz$, cho đường thẳng $d\colon \dfrac{x-1}{4}=\dfrac{-y}{2}=\dfrac{z+2}{-6}$. Vectơ nào dưới đây là một vectơ chỉ phương của đường thẳng $d$?
}{%
\begin{tasks}(4)
    \task $\vec{u}_1=(2;-1;3)$
    \task $\vec{u}_2=(4;2;-6)$
    \task $\vec{u}_3=(-2;1;3)$
    \task $\vec{u}_4=(1;0;2)$
\end{tasks}
}

\questionLinesManual{0}{4}{%
Trong không gian $Oxyz$, cho đường thẳng $d\colon \begin{cases} x=1+t \\ y=4 \\ z=3-2t \end{cases}$. Vectơ nào dưới đây là một vectơ chỉ phương của đường thẳng $d$?
}{%
\begin{tasks}(4)
    \task $\vec{u}_1=(1;4;3)$
    \task $\vec{u}_2=(1;4;-2)$
    \task $\vec{u}_3=(1;0;-2)$
    \task $\vec{u}_4=(1;0;2)$
\end{tasks}
}

\questionLinesManual{0}{5}{%
Trong không gian $Oxyz$, cho hai điểm $A(1;1;0)$ và $B(0;1;2)$. Vectơ nào dưới đây là một vectơ chỉ phương của đường thẳng $AB$?
}{%
\begin{tasks}(4)
    \task $\vec{u}_1=(-1;0;2)$
    \task $\vec{u}_2=(1;2;2)$
    \task $\vec{u}_3=(-1;1;2)$
    \task $\vec{u}_4=(-1;0;-2)$
\end{tasks}
}

\questionLinesManual{0}{6}{%
Trong không gian $Oxyz$, cho đường thẳng $d\colon \dfrac{x-2}{4}=\dfrac{y-1}{-2}=\dfrac{z+3}{1}$. Điểm nào dưới đây thuộc đường thẳng $d$?
}{%
\begin{tasks}(4)
    \task $Q(4;-2;1)$
    \task $N(4;2;1)$
    \task $P(2;1;-3)$
    \task $M(2;1;3)$
\end{tasks}
}

\questionLinesManual{0}{7}{%
Trong không gian $Oxyz$, cho đường thẳng $d\colon \dfrac{x-4}{2}=\dfrac{y-2}{-5}=\dfrac{z+1}{1}$. Điểm nào sau đây thuộc $d$?
}{%
\begin{tasks}(4)
    \task $N(4;2;-1)$
    \task $Q(2;5;1)$
    \task $M(4;2;1)$
    \task $P(2;-5;1)$
\end{tasks}
}

\questionLinesManual{0}{8}{%
Trong không gian $Oxyz$, cho đường thẳng $d\colon \begin{cases} x=1+2t \\ y=3-t \\ z=1-t \end{cases}$. Điểm nào dưới đây không thuộc đường thẳng $d$?
}{%
\begin{tasks}(4)
    \task $P(0;1;2)$
    \task $A(1;3;1)$
    \task $Q(3;2;0)$
    \task $N(5;1;-1)$
\end{tasks}
}

\questionLinesManual{0}{9}{%
Trong không gian $Oxyz$, cho đường thẳng $d\colon \begin{cases} x=4-2t \\ y=-3+t \\ z=1-t \end{cases}$. Giao điểm của $d$ với mặt phẳng tọa độ $(Oxy)$ có tọa độ là:
}{%
\begin{tasks}(4)
    \task $(4;-3;0)$
    \task $(2;-2;0)$
    \task $(0;-1;-1)$
    \task $(-2;0;-2)$
\end{tasks}
}

\questionLinesManual{0}{10}{%
Trong không gian $Oxyz$, cho đường thẳng $d\colon \dfrac{x-12}{4}=\dfrac{y-9}{3}=\dfrac{z-1}{1}$ và mặt phẳng $(P)\colon 3x+5y-z-2=0$. Tọa độ giao điểm của $d$ và $(P)$ là:
}{%
\begin{tasks}(4)
    \task $(1;0;1)$
    \task $(0;0;-2)$
    \task $(1;1;6)$
    \task $(0;0;2)$
\end{tasks}
}

% =========================================================================
% PHẦN 2: THÔNG HIỂU - VIẾT PHƯƠNG TRÌNH ĐƯỜNG THẲNG CƠ BẢN (CÂU 11 - CÂU 23)
% =========================================================================

\questionLinesManual{0}{11}{%
Trong không gian $Oxyz$, phương trình tham số của đường thẳng đi qua điểm $M(2;0;-1)$ và có vectơ chỉ phương $\vec{a}=(2;-3;1)$ là:
}{%
\begin{tasks}(2)
    \task $\begin{cases} x=4+2t \\ y=-6 \\ z=2-t \end{cases}$
    \task $\begin{cases} x=-2+2t \\ y=-3t \\ z=1+t \end{cases}$
    \task $\begin{cases} x=2+2t \\ y=-3t \\ z=-1+t \end{cases}$
    \task $\begin{cases} x=-2+4t \\ y=-6t \\ z=1+2t \end{cases}$
\end{tasks}
}

\questionLinesManual{0}{12}{%
Trong không gian $Oxyz$, phương trình đường thẳng đi qua điểm $A(1;2;3)$ và có vectơ chỉ phương $\vec{u}=(2;-1;-2)$ là:
}{%
\begin{tasks}(2)
    \task $\dfrac{x-2}{1}=\dfrac{y+1}{2}=\dfrac{z+2}{3}$
    \task $\dfrac{x+1}{2}=\dfrac{y+2}{-1}=\dfrac{z+3}{-2}$
    \task $\dfrac{x+2}{1}=\dfrac{y-1}{2}=\dfrac{z-2}{3}$
    \task $\dfrac{x-1}{2}=\dfrac{y-2}{-1}=\dfrac{z-3}{-2}$
\end{tasks}
}

\questionLinesManual{6}{13}{%
Trong không gian $Oxyz$, phương trình tham số của đường thẳng đi qua điểm $M(1;-2;4)$ và song song với trục hoành $Ox$ là:
}{%
\begin{tasks}(2)
    \task $\begin{cases} x=1+t \\ y=-2 \\ z=4 \end{cases}$
    \task $\begin{cases} x=1 \\ y=-2+t \\ z=4 \end{cases}$
    \task $\begin{cases} x=1 \\ y=-2 \\ z=4+t \end{cases}$
    \task $\begin{cases} x=t \\ y=-2t \\ z=4t \end{cases}$
\end{tasks}
}

\questionLinesManual{6}{14}{%
Trong không gian $Oxyz$, cho hai điểm $M(1;0;1)$ và $N(3;2;-1)$. Đường thẳng $MN$ có phương trình tham số là:
}{%
\begin{tasks}(2)
    \task $\begin{cases} x=1+2t \\ y=2t \\ z=1+t \end{cases}$
    \task $\begin{cases} x=1+t \\ y=t \\ z=1+t \end{cases}$
    \task $\begin{cases} x=1-t \\ y=t \\ z=1+t \end{cases}$
    \task $\begin{cases} x=1+t \\ y=t \\ z=1-t \end{cases}$
\end{tasks}
}

\questionLinesManual{6}{15}{%
Trong không gian $Oxyz$, cho hai điểm $E(-1;0;2)$ và $F(2;1;-5)$. Phương trình chính tắc của đường thẳng $EF$ là:
}{%
\begin{tasks}(2)
    \task $\dfrac{x-1}{3}=\dfrac{y}{1}=\dfrac{z+2}{-7}$
    \task $\dfrac{x+1}{3}=\dfrac{y}{1}=\dfrac{z-2}{-7}$
    \task $\dfrac{x-1}{1}=\dfrac{y}{1}=\dfrac{z+2}{-3}$
    \task $\dfrac{x+1}{1}=\dfrac{y}{1}=\dfrac{z-2}{3}$
\end{tasks}
}

\questionLinesManual{6}{16}{%
Trong không gian $Oxyz$, cho tam giác $ABC$ với $A(1;-3;4)$, $B(-2;-5;-7)$, $C(6;-3;-1)$. Phương trình đường trung tuyến $AM$ của tam giác là:
}{%
\begin{tasks}(2)
    \task $\begin{cases} x=1+t \\ y=-1-3t \\ z=-8-4t \end{cases}$
    \task $\begin{cases} x=1-3t \\ y=-3-2t \\ z=4-11t \end{cases}$
    \task $\begin{cases} x=1+t \\ y=-3-t \\ z=4-8t \end{cases}$
    \task $\begin{cases} x=1+3t \\ y=-3+4t \\ z=4-t \end{cases}$
\end{tasks}
}

\questionLinesManual{6}{17}{%
Trong không gian $Oxyz$, cho ba điểm $A(2;1;3)$, $B(1;0;1)$, $C(-1;1;2)$. Phương trình chính tắc của đường thẳng đi qua điểm $A$ và song song với đường thẳng $BC$ là:
}{%
\begin{tasks}(2)
    \task $\begin{cases} x=2-2t \\ y=1+t \\ z=3+t \end{cases}$
    \task $\dfrac{x+2}{-2}=\dfrac{y+1}{1}=\dfrac{z+3}{1}$
    \task $\dfrac{x-2}{-2}=\dfrac{y-1}{1}=\dfrac{z-3}{1}$
    \task $\dfrac{x+2}{2}=\dfrac{y-1}{1}=\dfrac{z-1}{3}$
\end{tasks}
}

\questionLinesManual{6}{18}{%
Trong không gian $Oxyz$, cho ba điểm $A(1;2;3)$, $B(1;1;1)$ và $C(3;4;0)$. Đường thẳng đi qua $A$ và song song với $BC$ có phương trình là:
}{%
\begin{tasks}(2)
    \task $\dfrac{x+1}{4}=\dfrac{y+2}{5}=\dfrac{z+3}{1}$
    \task $\dfrac{x-1}{2}=\dfrac{y-2}{3}=\dfrac{z-3}{-1}$
    \task $\dfrac{x-1}{4}=\dfrac{y-2}{5}=\dfrac{z-3}{1}$
    \task $\dfrac{x+1}{2}=\dfrac{y+2}{3}=\dfrac{z+3}{-1}$
\end{tasks}
}

\questionLinesManual{6}{19}{%
Trong không gian $Oxyz$, cho hai điểm $A(1;-2;-3)$, $B(-1;4;1)$ và đường thẳng $d\colon \dfrac{x+2}{1}=\dfrac{y-2}{-1}=\dfrac{z+3}{2}$. Phương trình đường thẳng đi qua trung điểm của $AB$ và song song với $d$ là:
}{%
\begin{tasks}(2)
    \task $\dfrac{x}{1}=\dfrac{y-1}{1}=\dfrac{z+1}{2}$
    \task $\dfrac{x}{1}=\dfrac{y-1}{-1}=\dfrac{z+1}{2}$
    \task $\dfrac{x-1}{1}=\dfrac{y-1}{-1}=\dfrac{z+1}{2}$
    \task $\dfrac{x}{1}=\dfrac{y-2}{-1}=\dfrac{z+2}{2}$
\end{tasks}
}

\questionLinesManual{6}{20}{%
Trong không gian $Oxyz$, cho điểm $M(1;2;-3)$ và mặt phẳng $(P)\colon 2x-y+3z-1=0$. Phương trình của đường thẳng đi qua điểm $M$ và vuông góc với $(P)$ là:
}{%
\begin{tasks}(2)
    \task $\begin{cases} x=2+t \\ y=-1+2t \\ z=3-3t \end{cases}$
    \task $\begin{cases} x=-1+2t \\ y=-2-t \\ z=3+3t \end{cases}$
    \task $\begin{cases} x=1+2t \\ y=2-t \\ z=-3+3t \end{cases}$
    \task $\begin{cases} x=1-2t \\ y=2-t \\ z=-3-3t \end{cases}$
\end{tasks}
}

\questionLinesManual{6}{21}{%
Trong không gian $Oxyz$, viết phương trình đường thẳng đi qua gốc tọa độ $O$ và vuông góc với mặt phẳng $(\alpha)\colon 2x-y-z-3=0$.
}{%
\begin{tasks}(2)
    \task $\begin{cases} x=-1+4t \\ y=1-2t \\ z=1-2t \end{cases}$
    \task $\begin{cases} x=2t \\ y=-t \\ z=-t \end{cases}$
    \task $\begin{cases} x=-2+2t \\ y=1+t \\ z=1-t \end{cases}$
    \task $\begin{cases} x=-2t \\ y=t \\ z=2t \end{cases}$
\end{tasks}
}

\questionLinesManual{6}{22}{%
Trong không gian $Oxyz$, phương trình chính tắc của đường thẳng đi qua điểm $M(1;-2;3)$ và vuông góc với mặt phẳng $(P)\colon x-2y+3z-1=0$ là:
}{%
\begin{tasks}(2)
    \task $\dfrac{x-1}{1}=\dfrac{y-2}{-2}=\dfrac{z-3}{3}$
    \task $\dfrac{x-1}{2}=\dfrac{y+2}{-2}=\dfrac{z-3}{3}$
    \task $\dfrac{x-1}{1}=\dfrac{y+2}{-2}=\dfrac{z-3}{3}$
    \task $\dfrac{x+1}{1}=\dfrac{y+2}{-2}=\dfrac{z-3}{3}$
\end{tasks}
}

\questionLinesManual{8}{23}{%
Trong không gian $Oxyz$, cho tam giác $ABC$ có ba đỉnh $A(3;0;0)$, $B(0;6;0)$, $C(0;0;6)$. Phương trình đường thẳng đi qua trực tâm $H$ của tam giác $ABC$ và vuông góc với mặt phẳng $(ABC)$ là:
}{%
\begin{tasks}(2)
    \task $\dfrac{x-2}{2}=\dfrac{y-1}{1}=\dfrac{z-1}{1}$
    \task $\dfrac{x-1}{2}=\dfrac{y-3}{1}=\dfrac{z-3}{1}$
    \task $\dfrac{x+1}{2}=\dfrac{y+2}{1}=\dfrac{z+3}{1}$
    \task $\dfrac{x-3}{2}=\dfrac{y-6}{1}=\dfrac{z-6}{1}$
\end{tasks}
}

% =========================================================================
% PHẦN 3: VẬN DỤNG VỪA SỨC - KẾT HỢP HAI ĐỐI TƯỢNG HÌNH HỌC (CÂU 24 - CÂU 42)
% =========================================================================

\questionLinesManual{6}{24}{%
Trong không gian $Oxyz$, cho hai mặt phẳng $(P)\colon 2x+y-z-1=0$ và $(Q)\colon x-2y+z-5=0$. Giao tuyến của $(P)$ và $(Q)$ có một vectơ chỉ phương là:
}{%
\begin{tasks}(4)
    \task $\vec{u}=(1;3;5)$
    \task $\vec{u}=(-1;3;-5)$
    \task $\vec{u}=(2;1;-1)$
    \task $\vec{u}=(1;-2;1)$
\end{tasks}
}

\questionLinesManual{8}{25}{%
Trong không gian $Oxyz$, cho hai mặt phẳng $x+z-5=0$ và $x-2y-z+3=0$. Giao tuyến của hai mặt phẳng đó có phương trình chính tắc là:
}{%
\begin{tasks}(2)
    \task $\dfrac{x+2}{1}=\dfrac{y+1}{3}=\dfrac{z}{-1}$
    \task $\dfrac{x+2}{1}=\dfrac{y+1}{2}=\dfrac{z}{-1}$
    \task $\dfrac{x-2}{1}=\dfrac{y-1}{1}=\dfrac{z-3}{-1}$
    \task $\dfrac{x-2}{1}=\dfrac{y-1}{2}=\dfrac{z-3}{-1}$
\end{tasks}
}

\questionLinesManual{6}{26}{%
Trong không gian $Oxyz$, cho điểm $A(1;2;-1)$ và hai mặt phẳng $(P)\colon 2x+3y=0$, $(Q)\colon 3x+4y=0$. Đường thẳng qua $A$ song song với cả $(P)$ và $(Q)$ có phương trình là:
}{%
\begin{tasks}(2)
    \task $\begin{cases} x=1+t \\ y=2+t \\ z=3+t \end{cases}$
    \task $\begin{cases} x=1 \\ y=2 \\ z=-1+t \end{cases}$
    \task $\begin{cases} x=t \\ y=2 \\ z=3+t \end{cases}$
    \task $\begin{cases} x=1 \\ y=t \\ z=3 \end{cases}$
\end{tasks}
}

\questionLinesManual{6}{27}{%
Trong không gian $Oxyz$, cho hai mặt phẳng $(\alpha)\colon x-2y+z-1=0$, $(\beta)\colon 2x+y-z=0$ và điểm $A(1;2;-1)$. Đường thẳng đi qua $A$ và song song với cả $(\alpha), (\beta)$ có phương trình là:
}{%
\begin{tasks}(2)
    \task $\dfrac{x-1}{-2}=\dfrac{y-2}{4}=\dfrac{z+1}{-2}$
    \task $\dfrac{x-1}{1}=\dfrac{y-2}{3}=\dfrac{z+1}{5}$
    \task $\dfrac{x-1}{1}=\dfrac{y-2}{-2}=\dfrac{z+1}{-1}$
    \task $\dfrac{x}{1}=\dfrac{y+2}{2}=\dfrac{z-3}{1}$
\end{tasks}
}

\questionLinesManual{8}{28}{%
Trong không gian $Oxyz$, cho điểm $A(1;2;1)$ và hai đường thẳng $d_1\colon \dfrac{x-1}{1}=\dfrac{y+1}{1}=\dfrac{z}{-1}$, $d_2\colon \dfrac{x+1}{2}=\dfrac{y-3}{1}=\dfrac{z-1}{2}$. Đường thẳng đi qua $A$ và vuông góc với cả $d_1, d_2$ có phương trình là:
}{%
\begin{tasks}(2)
    \task $\dfrac{x-1}{-3}=\dfrac{y-2}{4}=\dfrac{z-1}{2}$
    \task $\dfrac{x-1}{3}=\dfrac{y-2}{4}=\dfrac{z-1}{-1}$
    \task $\dfrac{x-1}{3}=\dfrac{y-2}{-4}=\dfrac{z-1}{-1}$
    \task $\dfrac{x+3}{-2}=\dfrac{y-4}{6}=\dfrac{z-1}{2}$
\end{tasks}
}

\questionLinesManual{8}{29}{%
Trong không gian $Oxyz$, cho điểm $I(2;5;3)$ và đường thẳng $d\colon \dfrac{x-1}{2}=\dfrac{y}{1}=\dfrac{z-2}{2}$. Đường thẳng $\Delta$ đi qua $I$ và vuông góc với hai đường thẳng $OI, d$ có phương trình là:
}{%
\begin{tasks}(2)
    \task $\dfrac{x+2}{7}=\dfrac{y+5}{-2}=\dfrac{z+3}{-8}$
    \task $\dfrac{x-2}{-8}=\dfrac{y-5}{7}=\dfrac{z-3}{-2}$
    \task $\dfrac{x+2}{7}=\dfrac{y+5}{2}=\dfrac{z+3}{-8}$
    \task $\dfrac{x-2}{7}=\dfrac{y-5}{2}=\dfrac{z-3}{-8}$
\end{tasks}
}

\questionLinesManual{8}{30}{%
Trong không gian $Oxyz$, cho điểm $A(1;2;3)$ và đường thẳng $d\colon \dfrac{x-3}{2}=\dfrac{y-1}{1}=\dfrac{z+7}{-2}$. Đường thẳng đi qua $A$, vuông góc với $d$ và cắt trục $Ox$ có phương trình là:
}{%
\begin{tasks}(2)
    \task $\begin{cases} x=-1+2t \\ y=-2t \\ z=t \end{cases}$
    \task $\begin{cases} x=1+t \\ y=2+2t \\ z=3+3t \end{cases}$
    \task $\begin{cases} x=1+t \\ y=2+2t \\ z=3+2t \end{cases}$
    \task $\begin{cases} x=-1+2t \\ y=2t \\ z=3t \end{cases}$
\end{tasks}
}

\questionLinesManual{8}{31}{%
Trong không gian $Oxyz$, cho điểm $A(2;1;3)$ và đường thẳng $d\colon \dfrac{x+1}{1}=\dfrac{y-1}{-2}=\dfrac{z-2}{2}$. Đường thẳng đi qua $A$, vuông góc với $d$ và cắt trục $Oy$ có phương trình là:
}{%
\begin{tasks}(2)
    \task $\begin{cases} x=2t \\ y=-3+4t \\ z=3t \end{cases}$
    \task $\begin{cases} x=2+2t \\ y=1+t \\ z=3+3t \end{cases}$
    \task $\begin{cases} x=2+2t \\ y=1+3t \\ z=3+2t \end{cases}$
    \task $\begin{cases} x=2t \\ y=-3+3t \\ z=2t \end{cases}$
\end{tasks}
}

\questionLinesManual{0}{32}{%
Trong không gian $Oxyz$, cho điểm $A(1;-1;3)$ và hai đường thẳng $d_1\colon \dfrac{x-3}{3}=\dfrac{y+2}{3}=\dfrac{z-1}{-1}$, $d_2\colon \dfrac{x-2}{1}=\dfrac{y+1}{-1}=\dfrac{z-1}{1}$. Đường thẳng đi qua $A$, vuông góc với $d_1$ và cắt $d_2$ có phương trình là:
}{%
\begin{tasks}(2)
    \task $\dfrac{x-1}{5}=\dfrac{y+1}{-4}=\dfrac{z-3}{2}$
    \task $\dfrac{x-1}{3}=\dfrac{y+1}{-2}=\dfrac{z-3}{3}$
    \task $\dfrac{x-1}{6}=\dfrac{y+1}{-5}=\dfrac{z-3}{3}$
    \task $\dfrac{x-1}{2}=\dfrac{y+1}{-1}=\dfrac{z-3}{3}$
\end{tasks}
}
\drawdashlinesfull

\questionLinesManual{8}{33}{%
Trong không gian $Oxyz$, cho điểm $M(1;0;1)$ và đường thẳng $d\colon \dfrac{x-1}{1}=\dfrac{y-2}{2}=\dfrac{z-3}{3}$. Đường thẳng đi qua $M$, vuông góc với $d$ và cắt trục $Oz$ có phương trình là:
}{%
\begin{tasks}(2)
    \task $\begin{cases} x=1-3t \\ y=0 \\ z=1+t \end{cases}$
    \task $\begin{cases} x=1-3t \\ y=0 \\ z=1-t \end{cases}$
    \task $\begin{cases} x=1-3t \\ y=t \\ z=1+t \end{cases}$
    \task $\begin{cases} x=1+3t \\ y=0 \\ z=1+t \end{cases}$
\end{tasks}
}

\questionLinesManual{8}{34}{%
Trong không gian $Oxyz$, cho đường thẳng $d\colon \dfrac{x-3}{1}=\dfrac{y-3}{3}=\dfrac{z}{2}$, mặt phẳng $(\alpha)\colon x+y-z+3=0$ và điểm $A(1;2;-1)$. Đường thẳng $\Delta$ đi qua $A$, cắt $d$ và song song với $(\alpha)$ có phương trình là:
}{%
\begin{tasks}(2)
    \task $\dfrac{x-1}{1}=\dfrac{y-2}{-2}=\dfrac{z+1}{-1}$
    \task $\dfrac{x-1}{1}=\dfrac{y-2}{2}=\dfrac{z+1}{1}$
    \task $\dfrac{x-1}{1}=\dfrac{y-2}{2}=\dfrac{z-1}{1}$
    \task $\dfrac{x-1}{-1}=\dfrac{y-2}{-2}=\dfrac{z+1}{1}$
\end{tasks}
}

\questionLinesManual{8}{35}{%
Trong không gian $Oxyz$, cho đường thẳng $d\colon \dfrac{x}{2}=\dfrac{y-3}{1}=\dfrac{z-2}{-3}$ và mặt phẳng $(P)\colon x-y+2z-6=0$. Đường thẳng nằm trong $(P)$, cắt và vuông góc với $d$ có phương trình là:
}{%
\begin{tasks}(2)
    \task $\dfrac{x+2}{1}=\dfrac{y-2}{7}=\dfrac{z-5}{3}$
    \task $\dfrac{x-2}{1}=\dfrac{y-4}{7}=\dfrac{z+1}{3}$
    \task $\dfrac{x-2}{1}=\dfrac{y+2}{7}=\dfrac{z+5}{3}$
    \task $\dfrac{x+2}{1}=\dfrac{y+4}{7}=\dfrac{z-1}{3}$
\end{tasks}
}

\questionLinesManual{0}{36}{%
Trong không gian $Oxyz$, cho mặt phẳng $(\alpha)\colon y+2z=0$ và hai đường thẳng $d_1\colon \begin{cases} x=1-t \\ y=t \\ z=4t \end{cases}$, $d_2\colon \begin{cases} x=2-t' \\ y=4+2t' \\ z=1 \end{cases}$. Đường thẳng $\Delta$ nằm trong mặt phẳng $(\alpha)$ và cắt cả hai đường thẳng $d_1, d_2$ có phương trình là:
}{%
\begin{tasks}(2)
    \task $\dfrac{x-1}{7}=\dfrac{y}{8}=\dfrac{z}{-4}$
    \task $\dfrac{x+1}{7}=\dfrac{y}{-8}=\dfrac{z}{4}$
    \task $\dfrac{x-1}{4}=\dfrac{y}{-2}=\dfrac{z}{1}$
    \task $\dfrac{x-1}{7}=\dfrac{y}{8}=\dfrac{z}{4}$
\end{tasks}
}
\drawdashlinesfull

\questionLinesManual{8}{37}{%
Trong không gian $Oxyz$, xét vị trí tương đối và tìm tọa độ giao điểm (nếu có) của hai đường thẳng $d_1\colon \begin{cases} x=1+2t \\ y=-1+3t \\ z=5+t \end{cases}$ và $d_2\colon \dfrac{x+1}{3}=\dfrac{y+4}{2}=\dfrac{z-4}{4}$. Khẳng định nào sau đây đúng?
}{%
\begin{tasks}(2)
    \task $d_1$ cắt $d_2$ tại $M(-1;-4;4)$
    \task $d_1$ và $d_2$ chéo nhau
    \task $d_1$ song song với $d_2$
    \task $d_1$ cắt $d_2$ tại $N(1;-1;5)$
\end{tasks}
}

\questionLinesManual{6}{38}{%
Trong không gian $Oxyz$, cho hai đường thẳng $d_1\colon \dfrac{x-1}{1}=\dfrac{y+2}{2}=\dfrac{z}{3}$ và $d_2\colon \dfrac{x+2}{2}=\dfrac{y-1}{1}=\dfrac{z-1}{1}$. Vị trí tương đối của $d_1$ và $d_2$ là:
}{%
\begin{tasks}(4)
    \task Chéo nhau
    \task Trùng nhau
    \task Song song nhau
    \task Cắt nhau
\end{tasks}
}

\questionLinesManual{6}{39}{%
Trong không gian $Oxyz$, cho hai đường thẳng $d_1\colon \dfrac{x-1}{2}=\dfrac{y+1}{-1}=\dfrac{z-2}{3}$ và $d_2\colon \begin{cases} x=2-4t \\ y=-2+2t \\ z=1-6t \end{cases}$. Vị trí tương đối của $d_1$ và $d_2$ là:
}{%
\begin{tasks}(4)
    \task Song song nhau
    \task Cắt nhau
    \task Chéo nhau
    \task Trùng nhau
\end{tasks}
}

\questionLinesManual{6}{40}{%
Trong không gian $Oxyz$, cho hai đường thẳng $d_1\colon \dfrac{x-1}{2}=\dfrac{y+1}{-1}=\dfrac{z-2}{3}$ và $d_2\colon \begin{cases} x=3-4t \\ y=-2+2t \\ z=5-6t \end{cases}$. Vị trí tương đối của $d_1$ và $d_2$ là:
}{%
\begin{tasks}(4)
    \task Trùng nhau
    \task Cắt nhau
    \task Song song nhau
    \task Chéo nhau
\end{tasks}
}

\questionLinesManual{0}{41}{%
Trong không gian $Oxyz$, cho điểm $A(1;0;-2)$ và đường thẳng $\Delta\colon \dfrac{x-1}{2}=\dfrac{y+1}{-1}=\dfrac{z}{1}$. Viết phương trình tham số của đường thẳng $d$ đi qua $A$, vuông góc và cắt $\Delta$.
}{%
\begin{tasks}(2)
    \task $\begin{cases} x=1+t \\ y=-t \\ z=-2-3t \end{cases}$
    \task $\begin{cases} x=1+2t \\ y=t \\ z=-2-3t \end{cases}$
    \task $\begin{cases} x=1+2t \\ y=-t \\ z=-2+t \end{cases}$
    \task $\begin{cases} x=1 \\ y=t \\ z=-2+t \end{cases}$
\end{tasks}
}
\drawdashlinesfull

\questionLinesManual{0}{42}{%
Trong không gian $Oxyz$, cho điểm $M(2;-1;3)$ và đường thẳng $\Delta\colon \begin{cases} x=1+2t \\ y=-2+t \\ z=1-t \end{cases}$. Phương trình chính tắc của đường thẳng $d$ đi qua điểm $M$, vuông góc và cắt $\Delta$ là:
}{%
\begin{tasks}(2)
    \task $\dfrac{x-2}{4}=\dfrac{y+1}{5}=\dfrac{z-3}{13}$
    \task $\dfrac{x-2}{1}=\dfrac{y+1}{-1}=\dfrac{z-3}{1}$
    \task $\dfrac{x-2}{1}=\dfrac{y+1}{2}=\dfrac{z-3}{1}$
    \task $\dfrac{x-2}{2}=\dfrac{y+1}{1}=\dfrac{z-3}{-1}$
\end{tasks}
}
\drawdashlinesfull

% =========================================================================
% PHẦN 4: VẬN DỤNG CAO - HÌNH CHIẾU, CẮT 2 ĐƯỜNG, VUÔNG GÓC CHUNG & KHOẢNG CÁCH (CÂU 43 - CÂU 60)
% =========================================================================

\questionLinesManual{0}{43}{%
Trong không gian $Oxyz$, tìm tọa độ hình chiếu vuông góc $H$ của gốc tọa độ $O$ lên mặt phẳng $(\alpha)\colon x+2y-z-6=0$.
}{%
\begin{tasks}(4)
    \task $H(1;2;-1)$
    \task $H(2;1;-2)$
    \task $H(-1;-2;1)$
    \task $H(1;1;-3)$
\end{tasks}
}
\drawdashlinesfull

\questionLinesManual{0}{44}{%
Trong không gian $Oxyz$, cho mặt phẳng $(\alpha)\colon 2x+y+2z-3=0$. Điểm $M'$ đối xứng với điểm $M(1;2;4)$ qua mặt phẳng $(\alpha)$ có tọa độ là:
}{%
\begin{tasks}(4)
    \task $(-3;0;0)$
    \task $(-1;1;2)$
    \task $(-1;-2;-4)$
    \task $(2;1;2)$
\end{tasks}
}
\drawdashlinesfull

\questionLinesManual{0}{45}{%
Trong không gian $Oxyz$, cho điểm $A(1;1;1)$ và đường thẳng $d\colon \begin{cases} x=6-4t \\ y=-2-t \\ z=-1+2t \end{cases}$. Tìm tọa độ hình chiếu vuông góc $A'$ của điểm $A$ trên đường thẳng $d$.
}{%
\begin{tasks}(4)
    \task $A'(2;3;1)$
    \task $A'(-2;3;1)$
    \task $A'(2;-3;1)$
    \task $A'(2;-3;-1)$
\end{tasks}
}
\drawdashlinesfull

\questionLinesManual{0}{46}{%
Trong không gian $Oxyz$, cho điểm $A(3;2;0)$ và đường thẳng $d\colon \dfrac{x+1}{1}=\dfrac{y+3}{2}=\dfrac{z+2}{2}$. Điểm đối xứng của $A$ qua đường thẳng $d$ có tọa độ là:
}{%
\begin{tasks}(4)
    \task $(-1;0;4)$
    \task $(7;1;-1)$
    \task $(2;1;-2)$
    \task $(0;2;-5)$
\end{tasks}
}
\drawdashlinesfull

\questionLinesManual{0}{47}{%
Trong không gian $Oxyz$, cho điểm $M(1;1;1)$ và hai đường thẳng $d_1\colon \dfrac{x-2}{1}=\dfrac{y+1}{2}=\dfrac{z}{-1}$, $d_2\colon \dfrac{x+1}{2}=\dfrac{y-1}{1}=\dfrac{z-1}{1}$. Viết phương trình tham số của đường thẳng $d$ đi qua điểm $M$ đồng thời cắt cả hai đường thẳng $d_1$ và $d_2$.
}{%
\begin{tasks}(2)
    \task $\begin{cases} x=1+t \\ y=1-t \\ z=1-t \end{cases}$
    \task $\begin{cases} x=1+2t \\ y=1+t \\ z=1+t \end{cases}$
    \task $\begin{cases} x=1+t \\ y=1-2t \\ z=1+3t \end{cases}$
    \task $\begin{cases} x=1 \\ y=1+t \\ z=1-t \end{cases}$
\end{tasks}
}
\drawdashlinesfull

\questionLinesManual{0}{48}{%
Trong không gian $Oxyz$, cho điểm $M(3;3;-2)$ và hai đường thẳng $d_1\colon \dfrac{x-1}{1}=\dfrac{y-2}{3}=\dfrac{z}{1}$, $d_2\colon \dfrac{x+1}{-1}=\dfrac{y-1}{2}=\dfrac{z-2}{4}$. Đường thẳng $d$ qua $M$ cắt $d_1, d_2$ lần lượt tại $A$ và $B$. Độ dài đoạn thẳng $AB$ bằng:
}{%
\begin{tasks}(4)
    \task $AB=2$
    \task $AB=3$
    \task $AB=\sqrt{6}$
    \task $AB=\sqrt{5}$
\end{tasks}
}
\drawdashlinesfull

\questionLinesManual{0}{49}{%
Trong không gian $Oxyz$, cho mặt phẳng $(P)\colon 2x-y+z-10=0$, điểm $A(1;3;2)$ và đường thẳng $d\colon \begin{cases} x=-2+2t \\ y=1+t \\ z=1-t \end{cases}$. Tìm phương trình đường thẳng $\Delta$ cắt $(P)$ và $d$ lần lượt tại $M, N$ sao cho $A$ là trung điểm cạnh $MN$.
}{%
\begin{tasks}(2)
    \task $\dfrac{x-6}{7}=\dfrac{y-1}{-4}=\dfrac{z+3}{-1}$
    \task $\dfrac{x+6}{7}=\dfrac{y+1}{4}=\dfrac{z-3}{-1}$
    \task $\dfrac{x-6}{7}=\dfrac{y-1}{4}=\dfrac{z+3}{-1}$
    \task $\dfrac{x+6}{7}=\dfrac{y+1}{-4}=\dfrac{z-3}{-1}$
\end{tasks}
}
\drawdashlinesfull

\questionLinesManual{0}{50}{%
Trong không gian $Oxyz$, cho ba đường thẳng $d_1\colon \dfrac{x-3}{2}=\dfrac{y+1}{1}=\dfrac{z-2}{-2}$, $d_2\colon \dfrac{x+1}{3}=\dfrac{y}{-2}=\dfrac{z+4}{-1}$ và $d_3\colon \dfrac{x+3}{4}=\dfrac{y-2}{-1}=\dfrac{z}{6}$. Đường thẳng song song với $d_3$, cắt cả $d_1$ và $d_2$ có phương trình là:
}{%
\begin{tasks}(2)
    \task $\dfrac{x-3}{4}=\dfrac{y+1}{1}=\dfrac{z-2}{6}$
    \task $\dfrac{x-3}{-4}=\dfrac{y+1}{1}=\dfrac{z-2}{-6}$
    \task $\dfrac{x+1}{4}=\dfrac{y}{-1}=\dfrac{z-4}{6}$
    \task $\dfrac{x-1}{4}=\dfrac{y}{-1}=\dfrac{z+4}{6}$
\end{tasks}
}
\drawdashlinesfull

\questionLinesManual{0}{51}{%
Trong không gian $Oxyz$, cho đường thẳng $\Delta\colon \dfrac{x-1}{1}=\dfrac{y+2}{-1}=\dfrac{z}{2}$ và hai đường thẳng $d_1\colon \dfrac{x-2}{2}=\dfrac{y-1}{1}=\dfrac{z+1}{-1}$, $d_2\colon \dfrac{x}{1}=\dfrac{y+1}{2}=\dfrac{z-2}{1}$. Phương trình tham số của đường thẳng $d$ song song với $\Delta$ đồng thời cắt cả $d_1$ và $d_2$ là:
}{%
\begin{tasks}(2)
    \task $\begin{cases} x=-\dfrac{5}{6}+t \\ y=-\dfrac{5}{12}-t \\ z=\dfrac{5}{12}+2t \end{cases}$
    \task $\begin{cases} x=t \\ y=-2-t \\ z=2t \end{cases}$
    \task $\begin{cases} x=2+t \\ y=1+t \\ z=-1-2t \end{cases}$
    \task $\begin{cases} x=-t \\ y=1-t \\ z=2+2t \end{cases}$
\end{tasks}
}
\drawdashlinesfull

\questionLinesManual{0}{52}{%
Trong không gian $Oxyz$, đường vuông góc chung của hai đường thẳng $d\colon \begin{cases} x=1+t \\ y=0 \\ z=-5+t \end{cases}$ và $d'\colon \begin{cases} x=0 \\ y=4-2t' \\ z=5+3t' \end{cases}$ có phương trình là:
}{%
\begin{tasks}(2)
    \task $\dfrac{x-4}{-1}=\dfrac{y}{3}=\dfrac{z+2}{1}$
    \task $\dfrac{x-4}{2}=\dfrac{y}{-3}=\dfrac{z-2}{-2}$
    \task $\dfrac{x+4}{-2}=\dfrac{y}{3}=\dfrac{z-2}{2}$
    \task $\dfrac{x-4}{-2}=\dfrac{y}{3}=\dfrac{z+2}{2}$
\end{tasks}
}
\drawdashlinesfull

\questionLinesManual{0}{53}{%
Trong không gian $Oxyz$, cho hai đường thẳng chéo nhau $d\colon \dfrac{x-3}{-4}=\dfrac{y+2}{1}=\dfrac{z+1}{1}$ và $d'\colon \dfrac{x}{-6}=\dfrac{y-1}{1}=\dfrac{z-2}{2}$. Phương trình đường thẳng vuông góc chung của $d$ và $d'$ là:
}{%
\begin{tasks}(2)
    \task $\dfrac{x+1}{1}=\dfrac{y+1}{2}=\dfrac{z}{2}$
    \task $\dfrac{x-1}{2}=\dfrac{y-1}{2}=\dfrac{z}{2}$
    \task $\dfrac{x+1}{1}=\dfrac{y-1}{2}=\dfrac{z}{2}$
    \task $\dfrac{x-1}{1}=\dfrac{y-1}{2}=\dfrac{z+1}{2}$
\end{tasks}
}
\drawdashlinesfull

\questionLinesManual{0}{54}{%
Trong không gian $Oxyz$, cho hai đường thẳng chéo nhau $d_1\colon \dfrac{x-1}{2}=\dfrac{y-2}{1}=\dfrac{z+1}{2}$ và $d_2\colon \dfrac{x+1}{1}=\dfrac{y-1}{1}=\dfrac{z-2}{-1}$. Tọa độ hai điểm $A\in d_1, B\in d_2$ sao cho $AB$ là đoạn vuông góc chung của $d_1$ và $d_2$ là:
}{%
\begin{tasks}(1)
    \task $A\left(\dfrac{22}{13};\dfrac{61}{26};-\dfrac{4}{13}\right)$ và $B\left(\dfrac{29}{26};\dfrac{81}{26};-\dfrac{3}{26}\right)$
    \task $A(1;2;-1)$ và $B(-1;1;2)$
    \task $A(5;4;3)$ và $B(0;2;1)$
    \task $A(-1;1;-3)$ và $B(2;4;-1)$
\end{tasks}
}
\drawdashlinesfull

\questionLinesManual{0}{55}{%
Trong không gian $Oxyz$, cho đường thẳng $d\colon \begin{cases} x=1+3t \\ y=-3 \\ z=5+4t \end{cases}$. Gọi $\Delta$ là đường thẳng đi qua điểm $A(1;-3;5)$ và có vectơ chỉ phương $\vec{u}=(1;2;-2)$. Đường phân giác của góc nhọn tạo bởi $d$ và $\Delta$ có phương trình là:
}{%
\begin{tasks}(2)
    \task $\begin{cases} x=-1+2t \\ y=2-5t \\ z=6+11t \end{cases}$
    \task $\begin{cases} x=-1+2t \\ y=2-5t \\ z=-6+11t \end{cases}$
    \task $\begin{cases} x=1+7t \\ y=-3+5t \\ z=5+t \end{cases}$
    \task $\begin{cases} x=1-t \\ y=-3 \\ z=5+7t \end{cases}$
\end{tasks}
}
\drawdashlinesfull

\questionLinesManual{8}{56}{%
Trong không gian $Oxyz$, cho điểm $A(2;1;-1)$ và đường thẳng $d\colon \dfrac{x-1}{2}=\dfrac{y+1}{1}=\dfrac{z}{-2}$. Khoảng cách từ điểm $A$ đến đường thẳng $d$ bằng:
}{%
\begin{tasks}(4)
    \task $\sqrt{2}$
    \task $3$
    \task $\sqrt{6}$
    \task $2$
\end{tasks}
}

\questionLinesManual{8}{57}{%
Trong không gian $Oxyz$, cho hai đường thẳng song song $d_1\colon \dfrac{x-1}{2}=\dfrac{y+2}{-1}=\dfrac{z-3}{1}$ và $d_2\colon \dfrac{x+2}{2}=\dfrac{y-1}{-1}=\dfrac{z+1}{1}$. Khoảng cách giữa hai đường thẳng $d_1$ và $d_2$ bằng:
}{%
\begin{tasks}(4)
    \task $d(d_1,d_2)=\sqrt{6}$
    \task $d(d_1,d_2)=3\sqrt{2}$
    \task $d(d_1,d_2)=\dfrac{\sqrt{210}}{6}$
    \task $d(d_1,d_2)=2\sqrt{3}$
\end{tasks}
}

\questionLinesManual{8}{58}{%
Trong không gian $Oxyz$, cho hai đường thẳng chéo nhau $d_1\colon \dfrac{x}{1}=\dfrac{y-3}{2}=\dfrac{z-2}{1}$ và $d_2\colon \dfrac{x-3}{1}=\dfrac{y+1}{-2}=\dfrac{z-2}{1}$. Khoảng cách giữa hai đường thẳng $d_1$ và $d_2$ bằng:
}{%
\begin{tasks}(4)
    \task $\dfrac{\sqrt{2}}{3}$
    \task $\dfrac{12}{5}$
    \task $\dfrac{3\sqrt{2}}{2}$
    \task $3$
\end{tasks}
}

\questionLinesManual{0}{59}{%
Trong không gian $Oxyz$, cho hai điểm $M(-2;-2;1)$, $A(1;2;-3)$ và đường thẳng $d\colon \dfrac{x+1}{2}=\dfrac{y-5}{2}=\dfrac{z}{-1}$. Tìm vectơ chỉ phương $\vec{u}$ của đường thẳng $\Delta$ đi qua $M$, vuông góc với đường thẳng $d$ đồng thời cách điểm $A$ một khoảng lớn nhất.
}{%
\begin{tasks}(4)
    \task $\vec{u}=(4;-5;-2)$
    \task $\vec{u}=(1;0;2)$
    \task $\vec{u}=(1;1;-4)$
    \task $\vec{u}=(8;-7;2)$
\end{tasks}
}
\drawdashlinesfull

\questionLinesManual{0}{60}{%
Trong không gian $Oxyz$, cho điểm $M(1;2;3)$, $A(1;0;0)$ và $B(0;0;3)$. Đường thẳng $\Delta$ đi qua $M$ và thỏa mãn tổng khoảng cách từ các điểm $A, B$ đến $\Delta$ là lớn nhất có phương trình là:
}{%
\begin{tasks}(2)
    \task $\Delta\colon \dfrac{x-1}{6}=\dfrac{y-2}{2}=\dfrac{z-3}{-3}$
    \task $\Delta\colon \dfrac{x-1}{6}=\dfrac{y-2}{-3}=\dfrac{z-3}{2}$
    \task $\Delta\colon \dfrac{x-1}{-3}=\dfrac{y-2}{6}=\dfrac{z-3}{2}$
    \task $\Delta\colon \dfrac{x-1}{2}=\dfrac{y-2}{-3}=\dfrac{z-3}{6}$
\end{tasks}
}
\drawdashlinesfull

% =========================================================================
% BẢNG ĐÁP ÁN BTVN 60 CÂU (ĐÃ KIỂM TRA ĐỘC LẬP - ẨN TRONG CODE)
% -------------------------------------------------------------------------
% 1.A  2.C  3.C  4.C  5.A  6.C  7.A  8.A  9.B  10.B
% 11.C  12.D  13.A  14.D  15.B  16.C  17.C  18.B  19.B  20.C
% 21.B  22.C  23.A  24.A  25.C  26.B  27.B  28.C  29.D  30.D
% 31.A  32.C  33.A  34.A  35.A  36.C  37.A  38.A  39.A  40.A
% 41.B  42.A  43.A  44.A  45.C  46.A  47.A  48.B  49.B  50.B
% 51.A  52.D  53.A  54.A  55.B  56.A  57.C  58.C  59.A  60.B
% =========================================================================

\end{document}